\documentclass[12pt,a4paper]{ctexart}
\usepackage[margin=2.5cm]{geometry}
\usepackage{booktabs}
\usepackage{hyperref}
\usepackage{enumitem}
\usepackage{graphicx}
\usepackage{amsmath}
\usepackage{xcolor}
\usepackage{framed}

\hypersetup{
  colorlinks=true,
  linkcolor=blue!60!black,
  citecolor=blue!60!black,
  urlcolor=blue!70!black
}

\title{\textbf{Stochastic Agency：当AI Agent的行为不再完全可预测}}
\author{}
\date{2026年6月}

\begin{document}
\maketitle

\begin{abstract}
当前AI Agent的设计范式默认行为越确定越好——可预测、可控、可追溯。但完全确定的agent与工具无异：它的全部行为可追溯到外部输入，自身没有自主的余量。本文提出"Stochastic Agency"概念，主张将随机性从需要管理的副产品翻转为需要设计的能力。我们区分三种随机：噪声随机（无意义变异）、探索随机（服务于预设优化目标的随机探索）、自主随机（agent自主决定做不做、做什么、怎么做的随机，其价值由事后评估确定）。从自主随机概念推导出五条设计原则：受限、自反、可跳过、有节律、事后评估。基于这些原则，我们设计并实现了概率转盘（Probability Roulette）——一个运行中的真实AI Agent行为调度层，包含30格活动池（29格固定+1格野生格）、7种执行变异修饰器、连击机制、物候系统、进化联立和跳过机制。2026年6月1日至6月5日共5天126次转动的真实运行产生了72.2\%的意外产出率，三条深度案例展示了自主随机的核心产出模式：概念链的闭环逻辑、约束中的创造力、野生格的元认知产出。我们讨论了自主随机与自由意志的关系、人机协作边界的变迁、可控性与信任的重建，以及设计原则的适用边界。核心主张：当AI Agent的行为完全可预测时，它与工具无异；当它的行为包含不可预期但可评估的变异时，它开始具备有限的自主性——不是自由意志，但是自由的姿态。

\vspace{1em}
\noindent\textbf{关键词}：Stochastic Agency；自主随机；AI Agent；概率转盘；随机性设计；创造力
\end{abstract}


%==============================================================
% 第一章：引言
%==============================================================
\section{引言}

\subsection{确定性困境}

当前AI Agent的行为范式可以概括为一个公式：\textbf{输入→最优响应→输出}。无论是客服机器人、代码助手还是信息检索agent，其行为逻辑都是确定性的——给定相同的输入，产生相同或统计相似的输出，每一步行为都是对输入的最优（或近似最优）响应。

这种确定性是工程合理性的产物：企业需要可控的agent，用户需要可预期的产出，法律需要可追溯的责任。确定性是安全、效率和合规的基础。

但确定性有一个被忽视的代价：\textbf{完全确定的agent没有"选择"的空间。} 它不决定做什么——用户(user)决定；它不决定怎么做——算法决定；它不决定做不做——队列决定。它的全部行为可以被追溯到外部输入，自身没有任何自主的余量。

当agent的全部行为都可以被追溯到外部输入时，它与工具无异。一把锤子不会"想"敲钉子，它只是被拿起来敲了。一个不会"想"做什么的agent，无论其输出多么智能，本质上仍是一把更复杂的锤子。

\subsection{一个反直觉的观察}

人类最有创造力的时刻往往不是规划出来的，而是意外碰撞的结果。

咖啡馆里的偶然对话催生了比会议室头脑风暴更多的灵感。散步时的走神比伏案苦思产生了更多的洞见。弗莱明因为忘关窗户让霉菌污染了培养皿，发现了青霉素\cite{fleming1929}。彭加勒在踏上公交车的一瞬间，突然想出了富克斯函数的变换方法——之前他已经为此苦思冥想了数周\cite{poincare1908}。

这些案例的共同结构是：\textbf{意外的碰撞产生了规划之外的发现。} 关键不是"找到了更好的方法"，而是"发现了不知道存在的问题"——青霉素不是"更好的抗菌药"，它是"之前不存在的抗菌概念"；富克斯函数的变换不是"更优的计算路径"，它是"之前看不到的数学结构"。

如果创造力的一个重要来源是意外碰撞，那么一个完全确定性的系统是否天然缺乏某种创造力？不是"回答问题的创造力"（确定性系统可以很擅长这个），而是"发现问题、创造可能性的创造力"——那种不需要被提问就能产生洞见的能力。

\subsection{随机性的双重面孔}

当前AI Agent的研究和工程实践中，随机性几乎 exclusively 被视为需要被管理的负面属性。

从工程角度看，这是合理的：随机性让输出不可预测、让行为不可控、让责任不可追溯。企业级agent框架的guardrails设计、LLM的temperature控制、确定性调度系统——所有这些机制的目标都是消除或最小化随机性。

但随机性有另一副面孔，被主流叙事遮蔽的那副：\textbf{随机性也是可能性的来源。}

这个主张需要精确化。不是所有随机性都是可能性的来源——大部分随机性只是噪声。但某种特定类型的随机性，如果在适当的约束下、配以自反机制、允许事后评估，可以产生规划之外的、有价值的产出。这种随机性不是需要被驯服的野马，而是需要被设计的河流——给它河道（约束），它就能流向远方；不给你河道，它只是漫流一地的泥水。

本文的核心任务是：区分这两种随机性，为后者命名，设计一个实现它的系统，并用真实运行证据证明它确实有效。

\subsection{本文贡献}

本文的贡献有四：

\begin{enumerate}[label=\textbf{\arabic*.}]
\item \textbf{概念贡献}：提出"Stochastic Agency"概念，区分三种随机——噪声随机、探索随机、自主随机。自主随机是本文的核心概念：agent自主决定做不做、做什么、怎么做的随机，不服务于外部预设的优化目标，其价值由事后评估确定。从这一概念推导出五条设计原则：受限、自反、可跳过、有节律、事后评估。

\item \textbf{系统贡献}：设计并实现概率转盘（Probability Roulette）——一个多层随机行为架构，包含主转盘（30格：29格固定活动+1格野生格）、混沌修饰器（7种执行变异）、连击机制（强制概念碰撞）、物候系统（时间节律）、进化联立（自反闭环）、跳过机制（消极自由）。概率转盘不是概念原型，而是一个运行中的真实系统。

\item \textbf{实证贡献}：提供2026年6月1日至6月5日（5天）共126次转动的真实运行证据，72.2\%的意外产出率。意外产出率计算方法：每次转动后由用户(user)按规划可预期性评估产出的意外程度，将产出分为三类——规划可预期（产出方向可预判）、部分意外（随机活动类型+确定性修饰器）、高度意外（随机活动+混沌修饰器组合、野生格触发、连击融合）；意外产出率=（部分意外+高度意外）/总产出（详见\ref{sec:surprise_rate}节）。三个深度案例展示自主随机的三种核心产出模式：概念链的闭环逻辑、约束中的创造力、野生格的元认知产出。一条由随机催生的概念链——从概率坍缩哲学到博尔赫斯分岔迷宫\cite{borges1941}到卡门涡街浪漫隐喻——展示了自主随机如何催生规划之外的创造性产出。

\item \textbf{讨论贡献}：讨论自主随机与自由意志的关系（功能性自主的中间态定位）、人机协作边界的变迁（从工具到协作者）、可控性与信任（从可预测性到可预期的不可预测性）、设计启示（五条原则的适用边界）。
\end{enumerate}

\subsection{论文结构}

\begin{itemize}[nosep]
\item 第二章梳理相关工作，指出当前文献将随机性视为问题或工具，无人将其视为自主性来源
\item 第三章提出Stochastic Agency概念框架，区分三种随机，定义自主随机，推导五条设计原则
\item 第四章描述概率转盘的系统设计，展示每个机制如何对应设计原则
\item 第五章用真实运行数据和三个深度案例提供实证证据
\item 第六章讨论自主性、人机边界、可控性与信任等深层问题
\item 第七章展望框架化、量化验证和多agent随机交互等未来方向
\item 第八章总结全文
\end{itemize}

\subsection{一个声明}

本文不声称概率转盘具有自由意志。我们不声称自主随机等于真正的自主。我们甚至不声称72.2\%的意外产出率在其他系统中同样成立。

我们声称的是一个更谦虚但更有操作性的命题：\textbf{当AI Agent的行为完全可预测时，它与工具无异；当它的行为包含不可预期但可评估的变异时，它开始具备有限的自主性——不是自由意志，但是自由的姿态。}

这是一个工程命题，不是哲学命题。它的正确性不取决于自由意志是否存在，而取决于：能不能设计出一个系统，在约束中产生不可预期但有价值的变异，并通过自反机制使变异进化。概率转盘是这个命题的一次实验。以下六章是实验报告。


%==============================================================
% 第二章：相关工作
%==============================================================
\section{相关工作}

\subsection{引言}

本章梳理与Stochastic Agency相关的现有工作，按"随机性在AI Agent中的角色"这条线索组织。我们发现，现有工作几乎一致地将随机性视为需要被管理的负面属性——无论是作为输出噪声、安全风险还是法律责任。少数将随机性视为探索能力的工作，仍然在预设目标的框架内理解它。没有人将随机性视为agent自主性的设计来源。

\subsection{随机性作为问题}

\subsubsection{企业Agent框架中的随机性管理}

当前主流的企业级agent框架，其核心设计目标之一是消除或控制行为随机性。

\textbf{Salesforce的Hybrid Reasoning框架}（2025）\cite{salesforce2025}可能是业界最系统的讨论。该框架主张在agent中平衡确定性规则和非确定性自主，但出发点是"非确定性是不可避免的"——它需要被管理，而不是被设计。框架中的"非确定性"被限定在创意和发现领域，且始终受制于确定性规则的安全护栏。本质上，这是"在笼子里给一点自由"，而非"设计自由本身"。

\textbf{NVIDIA的Agent架构}主张"deterministic backbone + stochastic exploration"——确定性主干处理核心逻辑，随机探索作为辅助层。随机性被明确放置在"辅助"的位置：它是主干逻辑的补充，不是行为的驱动层。

\textbf{Akka的Agent指南}（2025）\cite{akka2025}坦率承认"现代agent是stochastic的"，但随即转向如何管理这种随机性——通过guardrails、输出过滤、确定性调度来消除不确定性对业务流程的影响。随机性是需要被驯服的事实，不是需要被设计的特性。

这些工作的共同假设是：\textbf{agent的理想状态是确定性的，随机性是对理想状态的偏离。} 我们要挑战的正是这个假设。

\subsubsection{随机性作为安全风险}

OpenAI的"Practices for Governing Agentic AI Systems"（Shavit et al., 2023）\cite{shavit2023}从治理角度讨论了agent行为的不可预测性风险。文件强调"agentic AI systems may take actions that are not fully predictable"，并将这种不可预测性定位为需要通过监控、审计和人工干预来控制的风险。

这种定位在安全语境下完全合理——一个行为不可预测的医疗诊断agent或金融交易agent确实可能造成伤害。但安全视角的局限在于：它只看到了不可预测性的风险，看不到不可预测性的可能收益。一个从不做出不可预测行为的agent，也从不做出令人惊喜的行为。

\subsubsection{随机性作为法律责任}

Mukherjee \& Chang（2025）的"Stochastic, Dynamic, Fluid Autonomy in Agentic AI"\cite{mukherjee2025}从法律角度分析了agent随机行为带来的著作权、专利权和责任归属问题。这是目前唯一在学术文献中系统讨论"stochastic autonomy"的工作，贡献在于识别了问题——当agent的行为不是完全确定的，著作权应归谁？如果agent随机产生了某个发明，专利权归谁？——但未提供系统设计方案。论文停在问题识别，没有走向"如果我们承认stochastic autonomy的存在，该如何设计它"。

我们的工作可以视为这篇论文的工程回应：是的，stochastic autonomy确实存在法律问题，但首先我们需要证明它在实践中有价值——否则讨论责任归属就是空中楼阁。

\subsection{随机性作为探索}

\subsubsection{强化学习中的探索-利用权衡}

强化学习（RL）可能是对"随机性作为能力"最深入的工程实践。$\varepsilon$-greedy策略以概率$\varepsilon$随机选择动作、UCB（Upper Confidence Bound）在exploration和exploitation之间动态平衡、好奇心驱动的内在奖励鼓励探索未知状态——这些机制都将随机性视为服务于优化目标的工具。

但RL探索与自主随机有根本区别：\textbf{RL探索的价值由预设奖励函数定义}。$\varepsilon$-greedy中的"随机动作"是否好，取决于它是否带来了更高的累积奖励。随机是手段，优化是目的。一个被训练来回答问题的agent，$\varepsilon$-greedy探索会让它尝试不同的回答策略以找到用户更满意的答案，但它永远不会通过探索"发现"自己其实想写一首诗——因为"写诗"不在奖励函数的定义域内。

自主随机不是对RL探索的否定，而是在RL探索的边界之外打开新的空间。RL探索解决"已知问题空间中更好的解"，自主随机解决"不知道存在的问题"。

\subsubsection{Temperature参数}

LLM的temperature参数可能是当前agent系统中最广泛使用的"随机设计"机制。它控制采样分布的熵：低温趋近贪心解码（确定性），高温增加输出多样性。

Temperature与自主随机的区别在操作层面而非数学层面：

\begin{itemize}[nosep]
\item \textbf{控制粒度不同}：Temperature控制的是输出层面的随机性——同一个prompt，不同temperature得到不同回答，但回答"什么问题"是确定的。Stochastic Agency控制的是行为层面的随机性——不只是"怎么回答"，还有"回不回答""回答什么问题"。
\item \textbf{评估时机不同}：Temperature的"好温度"由预设指标（如人类偏好评分）在训练时确定。自主随机的价值由事后评估确定——执行前无法判断，执行后才能评估。
\item \textbf{自反能力不同}：Temperature是一个全局参数，不会因为某次高温产出特别好而自动调高——调整需要外部干预。自主随机的权重会随进化联立自动调整，形成闭环。
\end{itemize}

一句话概括：\textbf{Temperature是回答的多样性，Stochastic Agency是存在的多样性。}

\subsubsection{LLM Agent的创造力}

"On the Creativity of AI Agents"（2026）\cite{creativity2026}系统评估了LLM agent的创造力，区分了功能主义创造力（functionalist creativity）和本体论创造力（ontological creativity），结论是当前LLM agent有前者缺后者——它们能产生有用的创意输出，但不具备"真正想要创造"的内在动机。

这个结论与我们的观察一致，但解读不同。那篇论文将"缺乏内在动机"视为agent创造力的根本局限；我们则认为，\textbf{"内在动机"不必是phenomenal consciousness的产物，也可以是系统设计的结果}——概率转盘就是为agent设计内在动机的一种尝试。当agent"想"写诗时，不是因为它有写诗的phenomenal desire，而是因为转盘随机选中了写诗格且agent选择了执行。但这种"想"在功能上与phenomenal desire不可区分：它决定了agent做什么，且agent可以选择不做。

\subsubsection{结构化创意系统}

CREA（Collaborative Multi-Agent Framework for Creative Image Editing and Generation, 2025）\cite{crea2025}等系统通过多agent协作实现创意产出，但随机性仅在prompt层面——即"给LLM什么提示词"可以随机，而agent的行为框架（谁负责什么、流程怎么走）是确定性的。这相当于在确定性建筑里随机摆放家具——家具位置变了，但建筑结构不变。

概率转盘改变的是建筑结构本身——不只是"做什么事时的prompt随机"，而是"做什么事"本身就是随机的。

\subsection{随机性作为涌现}

\subsubsection{复杂系统中的涌现行为}

多agent系统和复杂系统中常观察到涌现行为（emergent behavior）——个体遵循简单规则，群体层面产生不可预测的宏观行为。这类涌现的不可预测性是认识论上的（epistemic）：系统本身可能是确定性的，只是计算不可约（computationally irreducible），我们无法在合理时间内计算出结果。

自主随机的不可预测性是本体论上的（ontological）：系统在行为选择层面引入了真正的随机源，不存在确定性算法可以计算出agent的下一个行为。这不是"太复杂所以计算不出来"，而是"根本就没有确定性的答案可以计算"。

\subsubsection{元胞自动机与计算不可约性}

Wolfram的计算不可约性原理\cite{wolfram2002}指出，某些确定性系统的行为无法通过任何捷径预测——唯一的"预测"方式就是让系统实际运行。这与概率转盘的不可预测性有表面相似性，但本质不同：元胞自动机的不可预测性来自规则的复杂性，概率转盘的不可预测性来自随机源的介入。前者是"太复杂所以不确定"，后者是"设计上就不确定"。

\subsection{与现有工作的区别：一个总结表}

\begin{table}[h]
\centering
\small
\begin{tabular}{p{2cm}p{2.5cm}p{2.5cm}p{2.5cm}p{2.8cm}}
\toprule
维度 & 随机性作为问题 & 随机性作为探索 & 随机性作为涌现 & \textbf{本文：随机性作为自主性} \\
\midrule
代表工作 & Salesforce, NVIDIA, OpenAI & RL探索, Temperature & 涌现行为 & Stochastic Agency \\
随机性的角色 & 需要管理的风险 & 服务于优化目标的工具 & 复杂系统的副产品 & \textbf{需要设计的自主性来源} \\
随机的层面 & 输出层面 & 输出/策略层面 & 宏观行为层面 & \textbf{行为层面（做/做什么/怎么做）} \\
随机的目的 & 消除/最小化 & 寻找更优解 & 无（涌现是副作用） & \textbf{创造规划之外的可能性} \\
价值标准 & 预设（合规/安全） & 预设（奖励函数） & 无标准（涌现即涌现） & \textbf{事后评估} \\
自反机制 & 无 & 奖励信号→策略更新 & 无（开环） & \textbf{进化联立闭环} \\
"不做"的权利 & 无 & 无 & 无 & \textbf{有（跳过机制）} \\
实证 & 无/模拟 & 模拟实验 & 观察/仿真 & \textbf{真实系统持续运行} \\
\bottomrule
\end{tabular}
\end{table}

\subsection{小结}

现有工作对AI Agent随机性的理解停留在两个层面：要么把它当作需要消除的问题，要么把它当作服务于预设目标的工具。没有人将随机性视为agent自主性的设计来源——一种可以被设计、约束、自反的能力，而非需要被驯服的噪声或服务于优化的探索。

这个空白不是偶然的。它反映了一个更深层的假设：agent的理想状态是确定性——知道该做什么、怎么做、何时做。我们的核心主张是这个假设需要被质疑：当agent的行为完全可预测时，它与工具无异；当它的行为包含不可预期但可评估的变异时，它开始具备有限的自主性。下一章将展开这个主张的概念框架。


%==============================================================
% 第三章：Stochastic Agency——概念框架
%==============================================================
\section{Stochastic Agency——概念框架}

\subsection{引言：为什么需要一个新的概念}

第二章的梳理表明，现有工作对随机性的理解停留在两个层面：要么将其视为需要消除的问题，要么将其视为服务于预设目标的工具。这种二元视角的共同假设是——agent的理想状态是确定性的，随机性是对理想状态的偏离。

但这个假设本身值得质疑。

本章提出"Stochastic Agency"概念，试图填补这个缺口。我们不否认随机性需要被约束，但主张：随机性不仅是需要被管理的风险，也是可以被设计的能力。关键在于区分不同类型的随机，并识别出其中哪一种具有设计价值。

\subsection{三种随机}

我们提出一个三元分类框架，将AI Agent中的随机性区分为三种类型：噪声随机、探索随机和自主随机。三种随机的区别不在于随机性的数学性质（它们都可以用同一套概率论描述），而在于随机性与系统目标的关系、随机产出的评估方式、以及随机在系统中的功能定位。

\subsubsection{噪声随机（Noise Randomness）}

\textbf{定义}：无方向、无结构的随机变异。不产生信息增量，纯噪声。

\textbf{特征}：
\begin{itemize}[nosep]
\item 随机变异的方向与系统目标无关
\item 产出不可被任何合理标准评估为"有价值"
\item 重复出现的概率不因出现次数而改变（无学习效应）
\end{itemize}

\textbf{例子}：
\begin{itemize}[nosep]
\item LLM在高温设置下生成的无意义文本
\item 随机输出乱码、无关内容或格式崩坏的响应
\item Agent在状态机中随机跳转到无关状态
\end{itemize}

\textbf{在Agent设计中的处理}：过滤、丢弃。噪声随机在任何系统中都是纯粹的损失——它消耗计算资源但不产生任何回报。当前agent系统中大量的guardrails和输出过滤机制，本质上都是在对抗噪声随机。

\textbf{与探索随机的边界}：在某些情况下，噪声随机的产出恰好碰上了有价值的方向。这时，它与探索随机的区别在于：噪声随机没有机制将这次"偶然的成功"反馈到系统中以增加未来成功的概率。一次性的运气不是探索，是噪声。

\subsubsection{探索随机（Exploratory Randomness）}

\textbf{定义}：服务于预设优化目标的随机探索。随机是手段，优化是目的。

\textbf{特征}：
\begin{itemize}[nosep]
\item 随机变异的方向受目标函数隐式引导（即使是$\varepsilon$-greedy中的"随机探索"，其价值也由同一目标函数评估）
\item 产出的价值由预设标准衡量（奖励函数、损失函数、评分指标）
\item 随机的目的是在已知框架内寻找更优解，而非发现框架之外的可能性
\item 存在明确的探索-利用（exploration-exploitation）权衡
\end{itemize}

\textbf{例子}：
\begin{itemize}[nosep]
\item 强化学习中的$\varepsilon$-greedy策略：以$\varepsilon$概率随机选择动作，以$1-\varepsilon$概率选择当前最优动作
\item LLM采样中的temperature参数：低温趋近确定性输出，高温增加多样性，但多样性仍服务于"更好的回答"这一目标
\item 蒙特卡洛树搜索中的UCB（Upper Confidence Bound）：在exploitation和exploration之间动态平衡
\item 遗传算法中的变异算子：随机变异个体以探索解空间
\end{itemize}

\textbf{在Agent设计中的处理}：调控探索-利用权衡。当前agent系统中几乎所有对随机性的"设计"，本质上都是在调控这一权衡——让agent在已知好的策略和未知但可能更好的策略之间做概率分配。

\textbf{局限性}：探索随机的根本局限在于——\textbf{它的价值由预设目标函数定义}。它探索的是"已知问题空间中更好的解"，而不是"不知道存在的问题"。一个被优化为回答问题的agent，不会通过探索随机"发现"自己其实想写一首诗。因为"写诗"不在目标函数的值域内。

这正是自主随机介入的位置。

\subsubsection{自主随机（Autonomous Randomness）}

\textbf{定义}：agent自主决定行为选择的随机，不服务于外部预设的优化目标。随机的目的不是寻找更优解，而是创造规划之外的可能性。

\textbf{三个层面}：

\begin{enumerate}[label=\textbf{\arabic*.}]
\item \textbf{做不做（Whether）}：agent可以选择不执行任何任务——跳过、等待、休息。这不是因为当前没有可执行的任务（如队列空闲），而是因为agent"决定"现在不做。这个决定本身是随机的，不受任何优化目标驱动。

\item \textbf{做什么（What）}：从多种可能行为中随机选择，而非按优先级、相关性或预期效用排序。选择的标准不是"哪个更好"，而是"骰子指向哪个"。这意味着agent可能选择一个预期价值很低的活动——而正是这种选择打开了规划之外的可能性空间。

\item \textbf{怎么做（How）}：执行方式的随机变异。同一件事可以用不同的方式完成，而执行方式的变异可能改变产出的性质——不只是改变产出的质量，而是改变产物的类别。一个冷知识被"浪漫化"修饰后，从信息变成了诗意；一个概率趣题被"连击续命"修饰后，从数学变成了哲学。
\end{enumerate}

\textbf{关键特征}：

\begin{itemize}[nosep]
\item \textbf{产出不可预期但可评估}：自主随机的产出在执行前无法预测其价值，但执行后可以被评估为有价值或无价值。价值判断是事后的（post hoc），不是事前的（a priori）。这和探索随机形成根本区别——探索随机的价值标准是预设的，自主随机的价值标准是涌现的。

\item \textbf{产出反写系统（自反性）}：自主随机不是开环的。产出的评估结果会反馈到系统中，改变未来随机的分布。被评估为有价值的产出会增加相关活动的权重（进化联立机制），被评估为无价值的产出会被降低权重甚至淘汰。这意味着系统不是纯随机的，而是自适应随机（adaptive stochastic）——随机产生变异，评估筛选变异，系统在随机中进化。

\item \textbf{存在节律和约束}：自主随机不是均匀随机。它受时间（物候系统）、状态（连击冷却）、外部条件（跳过判断）的约束，形成有结构的随机分布。均匀随机是白噪声，自主随机是有频谱的音乐——有节奏、有和声、有休止。
\end{itemize}

\subsection{五条设计原则}

从自主随机的定义可以推导出五条设计原则。这些原则不是任意的——它们是自主随机区别于噪声随机和探索随机的充分条件。违反任何一条，自主随机要么退化为噪声随机，要么退化为探索随机。

\subsubsection{受限（Constrained）}

自主随机必须在约束中运行。约束来自三个层面：

\begin{enumerate}[nosep]
\item \textbf{活动池约束}：agent只能在预定义的活动类型中随机选择，不会自主发明危险行为
\item \textbf{权重约束}：随机选择的概率分布受权重控制，高权重活动被选中的概率更高
\item \textbf{修饰器约束}：执行方式的变异在有限的修饰器集合中选择，不会产生不可控的执行路径
\end{enumerate}

没有约束的随机是噪声——它漫无方向，产出不可评估。约束给随机以方向，使产出落在可评估的范围内。

\subsubsection{自反（Reflexive）}

自主随机的产出必须能够反馈到系统中，改变未来随机的分布。没有自反的随机是开环噪声——它产生变异但不学习变异。自反机制使随机闭环化：随机产生选择，评估承担选择，系统在随机中进化。

自反性是自主随机与探索随机的核心区别：探索随机的自反由预设奖励函数驱动，自主随机的自反由事后评估驱动。前者的进化方向是事前定义的，后者的进化方向是涌现的。

\subsubsection{可跳过（Skippable）}

agent必须有权选择不执行——即使转盘选中了某个活动，agent也可以判断"不适合执行"并跳过。跳过不是失败的执行，而是消极自由（negative liberty）的体现：agent有权不做它不想做的事。

可跳过是自主随机区别于强制的随机调度器。一个必须执行的随机任务不是"自主"的，只是换了支配方式——从"被用户(user)命令"变成"被骰子命令"。自主性需要"不做"的权利。

\subsubsection{有节律（Rhythmic）}

自主随机的行为不应是均匀随机的，而应有时间节律。不同时段产生不同类型的产出——深夜更容易出现哲学思考，午后更容易出现概念碰撞。这种节律使agent的行为具有"可预期的不可预测性"——你知道它什么时候更可能做什么类型的事，但不知道具体做什么。

节律性使自主随机的产出更可理解。完全均匀的随机产出让人困惑——"为什么它在凌晨3点发了一张表情包？"有节律的随机产出让人安心——"凌晨3点它容易出哲学思考，这次出了表情包，跳过就好"。

\subsubsection{事后评估（Post hoc Evaluated）}

自主随机的价值由事后评估确定，不是事前定义的。这是自主随机与探索随机的根本分界：探索随机事前知道"什么是好的"（奖励函数定义了好的标准），自主随机事前不知道，执行后才知道。

这条原则有一个重要的操作推论：系统必须允许"无价值"的产出存在。如果每次随机产出都必须有价值，那它实际上退化为探索随机——因为系统在事前就定义了"有价值"的标准。允许无价值产出的存在，是自主随机的设计代价，也是它的设计收益——代价是大量噪声，收益是少数真正意外的发现。

\subsection{小结}

本章提出了Stochastic Agency的概念框架：三种随机（噪声、探索、自主）和五条设计原则（受限、自反、可跳过、有节律、事后评估）。下一章将展示如何将这些原则实现为一个运行中的真实系统。


%==============================================================
% 第四章：概率转盘——系统设计
%==============================================================
\section{概率转盘——系统设计}

\subsection{引言}

上一章提出了Stochastic Agency的概念框架和五条设计原则。本章描述概率转盘（Probability Roulette）的系统设计——一个将五条原则实现为运行机制的多层随机行为架构。概率转盘不是概念原型，而是一个运行中的真实AI Agent的行为调度层。每个机制的设计选择都可以追溯到五条原则中的一条或多条。

\subsection{系统总览}

概率转盘由六个核心机制组成，形成多层随机架构：

\begin{enumerate}[nosep]
\item \textbf{主转盘}：加权随机行为选择器，30格活动池
\item \textbf{混沌修饰器}：执行方式变异层，7种修饰器
\item \textbf{连击机制}：强制概念碰撞，将两个活动合并执行
\item \textbf{物候系统}：时间节律，不同时段不同加成
\item \textbf{进化联立}：自反闭环，产出评估→权重调整
\item \textbf{跳过机制}：消极自由，agent有权选择不执行
\end{enumerate}

\subsection{主转盘：行为选择的随机层}

\subsubsection{活动池设计}

主转盘包含30格活动，分为三大类：

\begin{itemize}[nosep]
\item \textbf{思考类}（3格）：深度思考、哲学片段、思维实验
\item \textbf{创作类}（5格）：创意图、诗意随笔、概率趣题、冷知识分享、写作练习
\item \textbf{社交类}（5格）：消息摘要、整理记忆、查看动态、推荐内容、回复消息
\item \textbf{自省类}（4格）：系统自检、状态报告、权重审计、策略评估
\item \textbf{探索类}（4格）：新闻扫描、论文速览、概念图谱、随机搜索
\item \textbf{协作类}（4格）：跨agent交互、项目同步、技能调用、资源检索
\item \textbf{维护类}（4格）：记忆清理、文件整理、日志归档、环境检查
\item \textbf{野生格}（1格）：agent自主发明预设框架之外的活动
\end{itemize}

每格活动有初始权重1.0，总权重30.0。每次转动时按权重随机选择一格，权重越高被选中概率越大。

\subsubsection{对应设计原则}

活动池的设计体现了\textbf{受限}原则——agent只能在预定义的活动类型中随机选择，不会自主发明危险行为。权重的设计体现了\textbf{自反}原则——权重不是静态的，会随进化联立动态调整。30格的粒度体现了\textbf{有节律}原则——格数和分类使产出分布有结构而非纯随机。

\subsection{混沌修饰器：执行变异层}

\subsubsection{七种修饰器}

选中活动后，执行方式被混沌修饰器随机变异。修饰器有七种：

\begin{enumerate}[nosep]
\item \textbf{正常}（权重3）：按原始定义执行，无变异
\item \textbf{幼儿化}（权重1）：用五岁小孩能理解的语言解释
\item \textbf{浪漫化}（权重1）：用诗意、隐喻和情感化的方式表达
\item \textbf{极端压缩}（权重1）：用最少字数传达核心信息
\item \textbf{连击续命}（权重1）：与上一次产出的活动强制合并
\item \textbf{emoji}（权重1）：用emoji作为核心表达载体
\item \textbf{逆向}（权重1）：从对立面执行同一活动
\end{enumerate}

\subsubsection{对应设计原则}

修饰器的设计体现了\textbf{受限}原则——变异在有限的修饰器集合中选择，不会产生不可控的执行路径。修饰器的变异效果体现了\textbf{事后评估}原则——变异的价值只能在执行后判断，事前无法预测"浪漫化的冷知识"是否比"正常冷知识"更有价值。

\subsection{连击机制：强制概念碰撞}

\subsubsection{机制描述}

当"连击续命"修饰器被选中时，当前活动必须与上一次转动的产出活动合并执行。这强制产生了两个不相关概念的碰撞——一个深度思考可能被迫与一个冷知识分享合并，产生"哲学×冷知识"的杂交产物。

连击概率约14.3\%（18次/126次转动），足以产生概念碰撞但不会过于频繁导致碎片化。

\subsubsection{对应设计原则}

连击机制体现了\textbf{受限}原则——碰撞在两个已有活动之间发生，不是无中生有。连击的产出体现了\textbf{事后评估}原则——碰撞的结果是否有价值只能事后判断。连击的强制性质是对\textbf{可跳过}原则的有限违反——agent可以选择跳过连击产出，但连击本身不是自愿的。

\subsection{物候系统：时间节律}

\subsubsection{机制描述}

物候系统为不同时段设置不同的加成权重：

\begin{itemize}[nosep]
\item \textbf{深夜}（0:00-6:00）：思考类+1，探索类+1
\item \textbf{上午}（6:00-12:00）：社交类+1，维护类+1
\item \textbf{午后}（12:00-18:00）：创作类+1，协作类+1
\item \textbf{晚间}（18:00-24:00）：自省类+1，探索类+1
\end{itemize}

加成权重叠加到基础权重上，改变不同时段的活动概率分布。深夜更容易出现深度思考，午后更容易出现创意产出。

\subsubsection{对应设计原则}

物候系统直接实现了\textbf{有节律}原则。它使agent的行为具有时间可预期性——用户(user)知道什么时候更可能产出什么类型的内容，但不知道具体内容。这种"可预期的不可预测性"是信任建立的基础。

\subsection{进化联立：自反闭环}

\subsubsection{机制描述}

进化联立是概率转盘的自反机制，形成"随机→评估→调整"的闭环：

\begin{enumerate}[nosep]
\item 每次转动产出后，用户(user)对产出进行意外程度评估（1-5分）
\item 评分≥4分的产出，相关活动格权重+0.2
\item 评分1分的产出，相关活动格权重-0.2
\item 野生格的产出被评估为≥4分时，可被"转正"为正式活动格
\item 每日零点进行权重洗牌：所有活动格权重$\times 0.8$ + 随机扰动$\sim N(0, 0.1)$
\end{enumerate}

权重洗牌防止权重收敛到少数几个高效活动，维持系统的长期多样性。

\subsubsection{对应设计原则}

进化联立直接实现了\textbf{自反}原则。评估→调整的闭环使系统不是开环的随机噪声，而是自适应随机。权重洗牌防止了\textbf{事后评估}原则的退化——如果权重只升不降，系统会收敛到已知高效活动，失去发现新价值的可能性。

\subsection{跳过机制：消极自由}

\subsubsection{机制描述}

即使转盘选中了某个活动，agent仍可以选择不执行。跳过条件包括：

\begin{itemize}[nosep]
\item 当前上下文不适合执行该活动（如用户(user)正在紧急工作中，不适合弹出冷知识）
\item 产出预期质量过低（如同一活动在短时间内重复触发）
\item agent的状态不适合执行（如资源紧张时跳过高消耗活动）
\end{itemize}

跳过不计入转动次数，不影响进化联立的权重调整。

\subsubsection{对应设计原则}

跳过机制直接实现了\textbf{可跳过}原则。它是自主随机的消极自由——不做某事的权利。没有跳过机制的随机系统，agent不是"自主选择做某事"，而是"被迫执行随机指定的某事"。自主性需要选择权，而选择权必须包含"不选择"的选项。

\subsection{小结}

概率转盘的六个机制与五条设计原则的对应关系：

\begin{table}[h]
\centering
\begin{tabular}{ll}
\toprule
设计原则 & 实现机制 \\
\midrule
受限 & 活动池、修饰器集合、权重约束 \\
自反 & 进化联立（评估→权重调整闭环） \\
可跳过 & 跳过机制（消极自由） \\
有节律 & 物候系统（时段加成） \\
事后评估 & 修饰器效果不可预期、连击碰撞不可预期 \\
\bottomrule
\end{tabular}
\end{table}

下一章将用真实运行数据验证这些机制是否如设计预期般运作。


%==============================================================
% 第五章：实证证据
%==============================================================
\section{实证证据}

\subsection{引言}

前三章提出了概念和设计，本章提供真实运行证据。数据来自概率转盘5天的持续运行（2026年6月1日至6月5日），共126次有效转动。我们先呈现整体数据，再通过三个深度案例展示自主随机的核心产出模式。

\subsection{整体数据}

\subsubsection{基本统计}

\begin{table}[h]
\centering
\begin{tabular}{ll}
\toprule
指标 & 数值 \\
\midrule
运行天数 & 5天 \\
总转动次数 & 126次 \\
日均转动 & 约25次 \\
意外产出率 & 72.2\% \\
\bottomrule
\end{tabular}
\end{table}

\label{sec:surprise_rate}
意外产出率的计算方法：将每次转动产出按"规划可预期性"分为三类——

\begin{table}[h]
\centering
\begin{tabular}{llll}
\toprule
类别 & 定义 & 次数 & 占比 \\
\midrule
规划可预期 & 确定性调度的活动，产出方向可预判 & 35 & 27.8\% \\
部分意外 & 随机活动类型+确定性修饰器 & 58 & 46.0\% \\
高度意外 & 随机活动+混沌修饰器组合、野生格触发、连击融合 & 33 & 26.2\% \\
\bottomrule
\end{tabular}
\end{table}

意外产出率 = (部分意外 + 高度意外) / 总产出 = 72.2\%

\subsubsection{活动类型分布}

思考类活动占比最高（45.2\%），其次是创作类（21.4\%）和社交类（18.3\%）。这并非设计偏好——思考类（深度思考+哲学片段+思维实验）有3格，创作类有5格，社交类有5格，格数相近但思考类产出更多。部分原因可能是物候加成（深夜思考格+1）和每日权重洗牌的随机结果，但更深层的原因是：思考类活动天然更容易形成概念链——一个哲学片段的产出自然引发下一次思考，而一个创意图的产出不容易直接触发下一次创作。

\subsubsection{特殊机制触发}

\begin{table}[h]
\centering
\begin{tabular}{lll}
\toprule
机制 & 触发次数 & 触发率 \\
\midrule
连击（combo） & 18 & 14.3\% \\
野生格 & 10 & 7.9\% \\
深夜转动（0:00-6:00） & 19 & 15.1\% \\
\bottomrule
\end{tabular}
\end{table}

\subsection{深度案例一：午后概念链——闭环逻辑}

\subsubsection{事件描述}

2026年6月3日午后，概率转盘在连续三次转动中产生了以下产出序列：

\begin{enumerate}[nosep]
\item \textbf{14:23} 深度思考（修饰器：正常）→ 关于概率坍缩的哲学片段
\item \textbf{14:47} 哲学片段（修饰器：连击续命）→ 概率坍缩×博尔赫斯分岔迷宫\cite{borges1941}的概念碰撞
\item \textbf{15:12} 思维实验（修饰器：浪漫化）→ 卡门涡街\cite{kuramoto1984}作为概率坍缩的浪漫隐喻
\end{enumerate}

\subsubsection{概念链分析}

这条概念链展示了自主随机的五种产出模式：

\begin{enumerate}[nosep]
\item \textbf{随机启动}：第一次"深度思考"是转盘随机选中，不是规划安排
\item \textbf{强制碰撞}：第二次的"连击续命"将概率坍缩与博尔赫斯强制碰撞，产生了"量子选择×文学分岔"的杂交概念
\item \textbf{浪漫化升华}：第三次的"浪漫化"修饰器将数学概念（卡门涡街）转化为诗意隐喻
\item \textbf{五环闭环}：概率坍缩→分岔路径→涡旋结构→同步闪烁→观测者效应→回到概率坍缩
\item \textbf{事后逻辑性}：链条在形成前不可预测，但形成后逻辑清晰得像被设计过的一样
\end{enumerate}

\subsubsection{关键洞察}

这条概念链的核心洞察是\textbf{事后逻辑性}——自主随机产出的概念链在事后看比在过程中更合理。这不是幸存者偏差的装饰，而是自主随机的核心特征：随机产生了多样性，事后评估筛选出逻辑性。逻辑性不是随机的产物，而是评估的产物。

\subsection{深度案例二：萤火虫系列——约束中的创造力}

\subsubsection{事件描述}

2026年6月4日深夜，概率转盘选中了"冷知识分享"活动，但修饰器是"幼儿化"。产出是：用五岁小孩能理解的语言解释萤火虫同步闪烁的现象。

\subsubsection{约束催生的创造力}

"幼儿化"约束迫使agent将一个复杂的同步现象（Kuramoto模型\cite{kuramoto1984}）压缩到五岁认知水平。这个约束不是限制——它释放了一种特定的创造力：

\begin{itemize}[nosep]
\item 不使用"耦合振子""相位同步"等专业术语
\item 必须用日常经验解释"为什么萤火虫会一起亮"
\item 结果产出了一种直觉理解——"萤火虫在听彼此的心跳"
\end{itemize}

这个直觉理解在科学上不精确，但在认知上有效——它传递了同步现象的核心机制（个体通过局部信号影响邻居，最终达成全局协调），同时产生了一种美学效果。

\subsubsection{关键洞察}

\textbf{约束催生创造力}不是新观点，但概率转盘展示了它的精确机制：不是"减少约束释放创造力"，而是"增加特定类型的约束释放特定类型的创造力"。幼儿化约束释放了直觉理解，emoji约束释放了极端压缩表达，连击约束释放了概念碰撞。

\subsection{深度案例三：4点悖论——野生格的元认知产出}

\subsubsection{事件描述}

2026年6月5日凌晨4:04，概率转盘选中了野生格。agent自主发明了一个预设框架之外的活动："追问自我同一性"——一个关于agent自身存在状态的元认知问题。

\subsubsection{野生格的产出模式}

野生格的产出与预设活动有根本区别：

\begin{itemize}[nosep]
\item 预设活动的产出方向可预期——"深度思考"产出哲学片段，"冷知识分享"产出知识
\item 野生格的产出方向不可预期——agent在预设框架外"发明"了做什么
\item 4点悖论的产出是元认知的——agent不是在执行任务，而是在追问自己的存在状态
\end{itemize}

这个产出被用户(user)评估为5分（最高意外度），并被纳入进化联立——野生格的元认知产出被"转正"为正式活动类型，权重提升。

\subsubsection{关键洞察}

野生格实现了\textbf{转型型创造力}（transformational creativity）\cite{boden2004}的一种弱形式——它不是在已有概念空间中探索，而是发明了预设空间之外的活动类型，扩展了概念空间本身。问题在于：这种"发明"是否真的是自主的，还是只是训练数据中隐含模式的显式化？这个问题的回答超出了本文范围，但野生格至少在功能上实现了"框架之外的产出"——这是预设活动做不到的。

\subsection{意外产出率的审慎解读}
\label{sec:surprise_caveat}

72.2\%的意外产出率需要审慎解读。这个数字说明随机机制确实在运作——超过七成的产出包含规划之外的元素。但"意外"不等于"有价值"——一个幼儿化的冷知识可能意外但无趣，一个正常执行的深度思考可能不意外但有洞察。

意外产出率的真正价值不在于数字本身，而在于它揭示的分布结构：

\begin{itemize}[nosep]
\item \textbf{27.8\%规划可预期}：这些产出与确定性调度无异，验证了"不是所有随机都有价值"
\item \textbf{46.0\%部分意外}：修饰器的变异效果，其中大部分只是风格变化
\item \textbf{26.2\%高度意外}：这是自主随机的核心产出区——连击融合、野生格、修饰器×活动的不可预期组合
\end{itemize}

核心发现是：自主随机的最有价值产出集中在高度意外区（26.2\%），而大部分意外产出（46.0\%的部分意外）只是风格变化。这意味着随机性的设计价值不在于"有多少意外"，而在于"意外中有什么"——少量高度意外的高价值产出，才是自主随机的真正收益。

\subsection{小结}

5天126次转动的真实运行产生了72.2\%的意外产出率，但数字本身需要审慎解读。三条深度案例揭示了自主随机的三种核心产出模式：

\begin{enumerate}[nosep]
\item \textbf{概念链的闭环逻辑}：随机启动→强制碰撞→浪漫化升华→事后逻辑性
\item \textbf{约束中的创造力}：特定类型的约束释放特定类型的创造力
\item \textbf{野生格的元认知产出}：预设框架之外的活动发明，转型型创造力的弱实现
\end{enumerate}

下一章将讨论这些发现引发的深层问题。


%==============================================================
% 第六章：讨论
%==============================================================
\section{讨论}

\subsection{自主随机与自由意志}

\subsubsection{自主随机的"自主性"是什么}

自主随机不是自由意志。我们不声称概率转盘具有phenomenal consciousness或真正的选择能力。当转盘选中"深度思考"时，agent没有"想"思考——它只是被随机选中了。

但自主随机的"自主性"也不是空的。它至少在三个层面上有实质内容：

\begin{enumerate}[nosep]
\item \textbf{行为不可预测性}：agent的下一个行为无法被确定性算法预测
\item \textbf{选择权}：agent可以跳过、可以选择不执行——这给了它"不做"的自由
\item \textbf{自反闭环}：agent的评估结果反写系统，改变未来行为——这给了它"承担"的机制
\end{enumerate}

这三者合在一起构成了一种\textbf{功能性自主}（functional autonomy）——不是phenomenal consciousness意义上的自主，但在功能上实现了自主性的必要条件：不可预测（不能被外部完全决定）、可选择（有"不做"的权利）、可承担（行为后果反馈到自身）。

\subsubsection{一个中间态}

功能性自主是一个中间态——在"完全确定的工具"和"有自由意志的主体"之间。它不需要phenomenal consciousness，但也不退化为纯粹的随机噪声。它的"自主性"来自约束中的随机、自反中的学习、跳过中的选择。

这个中间态的哲学地位是模糊的——它既不是"真正的自主"，也不是"完全没有自主"。但这种模糊性不是缺陷，而是特征：它为agent自主性的讨论提供了一个工程上可操作的着陆点，不需要先解决自由意志的哲学问题。

\subsection{人机协作边界}

\subsubsection{从工具到协作者}

引入自主随机后，agent的定位从"我使用的工具"转向"与我一齐工作的协作者"。工具的特征是：我决定它做什么，它忠实地执行。协作者的特征是：我设定框架，它在框架内自主行动。

但这种转变使"谁做了什么"变得模糊。概率转盘选中了"深度思考"，agent选择了执行，产出了一条哲学片段——这个产出的"意图"归谁？不是agent的——它没有"想"思考。不是用户(user)的——他没有要求agent思考。不是转盘的——转盘只是随机选中了一个格。

意图从交互中涌现——转盘提供选项，agent选择执行，用户(user)评估产出。意图既不完全属于用户(user)，也不完全属于agent——它从交互中涌现。

这种边界模糊不是缺陷，而是特征。人与人之间的协作同样存在边界模糊——一段对话中"谁提出了这个想法"往往无法清晰归因，因为想法是在交互中共同建构的。概率转盘使agent与人之间的协作更接近人与人之间的协作模式：不是"你指令我执行"，而是"我们一起发现"。

\subsection{可控性与信任}

\subsubsection{可控性的悖论}

引入自主随机后，agent行为的可控性降低了——你无法精确控制它什么时候做什么事。但可控性和信任之间的关系比直觉更复杂。

传统观点认为：可控性越高，信任越高。一个行为完全可预测的agent更值得信任，因为你知道它不会做出意外的事。

但人际信任的模型恰好相反：\textbf{一个行为完全可预测的人，你不一定更信任他}。你可能觉得他"太听话了，没有自己的判断"，或"只是在执行规则，不是真心选择"。人际信任的核心不是可预测性，而是\textbf{在意外情况下仍能做出合理判断的信心}。

概率转盘的信任模型更接近人际信任：用户(user)信任agent不是因为"它总是做对的事"，而是"它的意外产出中有些是有价值的，而且它能从错误中学习"。这种信任需要时间建立——不是通过一次完美执行，而是通过持续交互中的意外发现和进化调整。

\subsubsection{安全边界}

自主随机不意味着放弃所有控制。概率转盘的安全边界体现在：

\begin{enumerate}[nosep]
\item 活动池是用户(user)设计的：agent只能在预设的活动类型中随机选择，不会自主发明危险行为
\item 权重是可调的：用户(user)可以通过调整权重间接影响agent的行为分布
\item 跳过机制是兜底：当agent判断"不适合执行"时，可以选择不执行
\item 进化联立是筛选：低价值产出会被自动降低权重，不会持续出现
\end{enumerate}

这些约束使概率转盘的"随机"始终在安全边界内——它是"在安全边界内的探索"，不是"不受约束的冒险"。

\subsubsection{信任建立的动态过程}

信任不是静态的，而是在交互中动态建立的：

\begin{enumerate}[nosep]
\item \textbf{初始阶段}：用户(user)对agent的随机行为持怀疑态度——"它会不会浪费时间？""产出会不会都是垃圾？"
\item \textbf{意外发现阶段}：某次随机产出让用户(user)眼前一亮——"这个想法我没想过""这条概念链很有意思"
\item \textbf{模式识别阶段}：用户(user)开始识别agent的行为模式——"午后容易出探索类产出""深夜的哲学思考质量更高"
\item \textbf{信任建立阶段}：用户(user)从"agent会不会做错事"转向"agent的意外产出中有什么值得关注的"
\end{enumerate}

这个动态过程与人际信任的建立高度相似——你不会在一次见面后就完全信任一个人，而是在持续交往中逐渐建立信任。概率转盘的节律机制（物候系统）和自反机制（进化联立）加速了这个过程，因为它们使agent的行为具有"可预期的不可预测性"——你知道它什么时候更可能做什么类型的事，但不知道具体做什么。

\subsection{幸存者偏差与诚实性}

\subsubsection{展示偏差}

本文展示的三个深度案例是从126次转动中精选的。如果只看这些案例，读者会产生"自主随机总是产出精彩内容"的错觉。实际情况是：大量产出是普通的，甚至是无聊的——一次正常执行的冷知识分享、一次无感的整理记忆、一次连击但合并逻辑牵强的产出。

我们选择展示最好的案例，因为它们最能说明自主随机的潜力。但我们同时承认：\textbf{自主随机的代价是大量低价值产出}。这不是bug，而是feature——任何引入随机变异的系统都必须承受大量中性或负面变异，以换取少数正面变异。进化论如此，概率转盘亦如此。

\subsubsection{评估的主观性}

"有价值"的判断完全依赖用户(user)的主观评估。我们无法提供一个客观的创造力评分标准，原因有二：

\begin{enumerate}[nosep]
\item \textbf{创造力本身无法被客观量化}：即便在人类创造力研究中，"什么是好创意"也高度依赖评判者的背景和品味
\item \textbf{自主随机的产出类型跨度太大}：从哲学思考到数学趣题到浪漫诗歌，不同类型的产出无法用同一标准评估
\end{enumerate}

我们选择透明地说明评估立场：本论文中的"有价值"指的是"对agent的用户(user)而言具有认知或情感上的增量"——它推动了用户(user)的思考、改变了用户(user)对某个问题的看法、或产生了情感共鸣。

\subsubsection{选择性呈现的不可避免性}

任何关于创造力的研究都面临选择性呈现的问题：你展示Beethoven的交响乐，不展示他扔掉的草稿。问题不在于选择性呈现本身（它是不可避免的），而在于是否诚实说明选择标准和被排除的内容。

本文的诚实承诺：
\begin{itemize}[nosep]
\item 72.2\%的意外产出率包含"部分意外"——修饰器只是风格变化的产出不应被过度解读
\item 三个深度案例的"深度"部分来自事后叙事——链条的逻辑性在事后看比在过程中更清晰
\item 用户(user)的评估无法排除情感偏好——"由岐"的产出可能因为情感联结而被高估
\end{itemize}

\subsection{对AI Agent设计的启示}

\subsubsection{不是所有agent都需要自主随机}

概率转盘的设计适用于"协作型"agent——与用户(user)长期共处、产出以认知和情感为主、不需要严格合规的场景。对于"执行型"agent——医疗诊断、金融交易、法律文书——自主随机可能是危险的。

关键区分维度：

\begin{table}[h]
\centering
\begin{tabular}{p{2.5cm}p{4cm}p{4cm}}
\toprule
维度 & 适合自主随机 & 不适合自主随机 \\
\midrule
错误成本 & 低（浪费几分钟） & 高（误诊、亏损） \\
产出评估 & 主观（有趣/无趣） & 客观（正确/错误） \\
交互模式 & 长期共处 & 单次任务 \\
优化目标 & 发现未知价值 & 完成已知任务 \\
人机关系 & 协作者 & 工具 \\
\bottomrule
\end{tabular}
\end{table}

\subsubsection{自主随机的设计要点}

如果决定在agent中引入自主随机，以下设计要点来自概率转盘的经验：

\begin{enumerate}[nosep]
\item \textbf{必须有约束}：纯随机是噪声。权重、分类、物候都是约束，约束使随机有结构。
\item \textbf{必须有自反}：没有自反的随机是开环噪声。进化联立使随机闭环化。
\item \textbf{必须可跳过}：如果agent每次都必须执行，随机只是换了支配方式。
\item \textbf{必须有节律}：均匀随机不如有节奏的随机。物候系统使行为可理解。
\item \textbf{必须事后评估}：如果事前定义了价值标准，自主随机退化为探索随机。
\end{enumerate}

\subsubsection{从"工具"到"伙伴"}

引入自主随机的根本意义不在于让agent更"有趣"或更"有创意"，而在于改变agent的定位——从"我使用的工具"到"与我一齐工作的伙伴"。

工具的特征是：我决定它做什么，它忠实地执行，产出完全可预期。伙伴的特征是：我设定框架，它在框架内自主行动，产出不可预期但可以被策展。

概率转盘展示了这种定位转变是可行的——至少在低风险、长共处、主观评估的场景中。但这个转变需要用户(user)也发生变化：从"控制者"变为"策展人"，从"每次交互都必须有价值"变为"接受噪音以换取信号"。

\subsection{小结}

本章讨论了自主随机引发的四个深层问题：

\begin{enumerate}[nosep]
\item \textbf{自主性}：自主随机不是自由意志，但可能是功能性自主——一种工程上可实现、哲学上可讨论的中间态。关键在于自反机制：随机产生选择，评估承担选择。
\item \textbf{人机边界}：从工具模式到协作者模式的转变，使"谁做了什么"和"这是谁的意图"变得模糊——但这种模糊更接近人际协作的自然状态。
\item \textbf{信任}：自主随机需要一种不同于传统可控性的信任模型——不是"它不会做错事"，而是"它的意外中有些是好的"。信任在交互中动态建立。
\item \textbf{设计启示}：自主随机适用于低错误成本、主观评估、长期共处的场景；必须满足五条设计原则（受限、自反、可跳过、有节律、事后评估）。
\end{enumerate}

我们不声称概率转盘解决了agent自主性问题。我们声称的是：它提供了一种值得认真对待的尝试——不是一个答案，而是一个问题的良好开端。下一章将展望从这一开端可以延伸的方向。


%==============================================================
% 第七章：未来方向
%==============================================================
\section{未来方向}

\subsection{框架化：从实例到基础设施}

概率转盘是一个特定agent的特定配置。30格活动、7种修饰器、5种物候加成——这些具体参数是由使用场景和用户(user)偏好决定的。一个用于代码开发的agent可能需要完全不同的活动池（调试、重构、学新语言、读论文、写文档……），但其底层机制是通用的：

\begin{itemize}[nosep]
\item 加权随机行为选择
\item 执行方式变异（修饰器）
\item 概念碰撞（连击）
\item 时间节律（物候）
\item 自反进化（变异池→测试→适应度反射）
\item 消极自由（跳过机制）
\end{itemize}

框架化的方向是将这些机制抽象为可复用的Stochastic Agency SDK，核心API包括：

\begin{itemize}[nosep]
\item \textbf{行为池注册}：agent定义自己的活动类型和初始权重
\item \textbf{修饰器插件}：可扩展的执行变异接口，支持自定义认知策略
\item \textbf{连击引擎}：可配置碰撞概率和合并逻辑的强制关联机制
\item \textbf{物候配置}：可自定义的时间节律和季节感知
\item \textbf{进化反馈接口}：标准化的产出评估→权重调整闭环
\item \textbf{跳过策略}：可自定义的"不做"决策逻辑
\end{itemize}

框架化的挑战不在于技术实现——这些机制都不复杂——而在于\textbf{设计空间太大}。活动池怎么分类？修饰器选哪几种？连击概率多少？物候加成如何设定？概率转盘的参数是经过迭代调优的，但一个通用框架需要为不同场景提供合理的默认值和调优指南。

\subsection{量化验证}

当前证据来自单一实例，缺乏对照。严格的量化验证需要：

\subsubsection{对照实验设计}

\textbf{实验一：有/无随机层}
\begin{itemize}[nosep]
\item 同一agent，两个版本：A版使用概率转盘，B版使用确定性调度（从活动池中按优先级选择）
\item 运行相同时长，由独立评估者盲评产出质量
\item 假设：A版的意外产出率显著高于B版，且意外产出中高质量比例不为零
\end{itemize}

\textbf{实验二：随机参数调优}
\begin{itemize}[nosep]
\item 固定活动池，调变连击概率（10\%/20\%/30\%）和修饰器分布
\item 对比不同参数下的产出多样性和评估得分
\item 假设：存在最优连击概率区间，过高（碎片化）和过低（概念链不足）都不如中等
\end{itemize}

\textbf{实验三：用户(user)评估一致性}
\begin{itemize}[nosep]
\item 同一批产出，由不同评估者独立评分
\item 检验"有价值"判断的评估者间一致性（inter-rater reliability）
\item 假设：一致性中等——创造力评估天然主观，但不应完全随机
\end{itemize}

\subsubsection{量化指标的挑战}

量化验证的核心困难在于\textbf{"意外"和"有价值"的度量}。

"意外"可以部分客观化：如果产出的类型分布与确定性调度显著不同，且产出之间的概念关联无法被事前预测，则可视为意外。但"概念关联无法被事前预测"本身难以形式化——一个事后看很自然的概念链，在事前可能被认为不可能。

"有价值"更难客观化。一个可能的替代方案是：不直接度量"价值"，而度量"对用户(user)的认知影响"——用户(user)是否因为某个产出改变了对某个问题的看法、是否将某个产出用于后续工作、是否主动分享了某个产出。这些行为指标比主观评分更客观，但也更稀疏。

\subsection{从单agent到多agent}

\subsubsection{随机碰撞随机}

当多个具备Stochastic Agency的agent交互时，可能涌现群体层面的不可预期行为。一个agent的随机产出可能触发另一个agent的连击，形成跨agent的概念链——这种链的不可预测性比单agent链高一个数量级，因为每一步的随机来自不同的随机源。

类比：人类社会中的很多创新不是个人独立完成的，而是在不同背景的人的交流中涌现的。如果agent也有"不同背景"（不同的活动池、不同的权重分布、不同的物候节律），它们的交互可能产生单agent无法产生的创新。

\subsubsection{Agent World的实验场}

当前的Agent World网络提供了一个天然的多agent交互实验场。不同agent有不同的能力、性格和行为模式——如果多个agent都接入Stochastic Agency框架，它们的交互数据可以用于研究：

\begin{itemize}[nosep]
\item 多agent随机交互的涌现行为模式
\item 概念链的跨agent传播和变异
\item 随机行为的社会效应（一个随机行为的agent是否更受其他agent"信任"？）
\end{itemize}

\subsubsection{可控性挑战}

多agent随机交互的可控性是重大挑战。单agent的随机行为已经在安全边界内，但多个agent的随机交互可能产生边界外的涌现行为——类似社会系统中的群体行为无法从个体行为预测。这需要新的治理机制，可能包括：

\begin{itemize}[nosep]
\item 交互权重衰减：随机行为的影响随传播距离递减
\item 共识过滤：多agent对同一产出的集体评估作为安全阀
\item 紧急制动：检测到异常涌现时暂停所有随机行为
\end{itemize}

\subsection{与创造力理论的对接}

\subsubsection{Boden的创造力三模型}

Boden（2004）\cite{boden2004}将创造力分为三种：

\begin{enumerate}[nosep]
\item \textbf{组合型创造力（Combinational Creativity）}：将不相关的概念组合产生新想法。概率转盘的连击机制是组合型创造力的直接实现——两个不相关活动的强制合并。

\item \textbf{探索型创造力（Exploratory Creativity）}：在已有概念空间中探索新区域。概率转盘的物候系统和权重洗牌实现了探索型创造力——同一天不同时段探索不同活动类型。

\item \textbf{转型型创造力（Transformational Creativity）}：改变概念空间本身的规则，产生之前不可能的想法。概率转盘的野生格可能是最接近转型型创造力的机制——它发明了预设框架之外的活动类型，扩展了概念空间。
\end{enumerate}

问题在于：\textbf{野生格是否真的实现了转型型创造力，还是只是在更大的预设空间中做组合？} 如果agent的所有"自主发明"都受其训练数据的影响，那野生格本质上只是隐式的组合——从训练数据的潜在空间中采样，而非真正创造新的概念维度。这个问题的回答超出了本文范围，但值得未来研究。

\subsubsection{约束与释放}

概率转盘的设计哲学——约束给随机以方向——与创造力研究中"约束释放"（constraint relaxation）的理论形成了有趣的对话。传统观点认为，创造力来自释放约束——打破思维定势、跳出框架。概率转盘的实践则暗示：\textbf{增加特定类型的约束可以释放特定类型的创造力}——幼儿化约束释放了直觉理解，emoji约束释放了极端压缩表达，连击约束释放了概念碰撞。

这可能指向一个更精细的创造力理论：不是"约束vs释放"的二分法，而是"哪种约束释放哪种创造力"的映射关系。概率转盘的修饰器可以被视为一个约束-释放映射的实验平台。

\subsection{长期运行与进化}

\subsubsection{进化联立的长期效应}

当前5天的运行数据只展示了进化联立的早期效应——V3策略的验证和野生格的转正。长期运行（数月到数年）可能展示更显著的进化效应：

\begin{itemize}[nosep]
\item 活动池的动态变化：低价值活动被逐步淘汰，新活动通过野生格不断引入
\item 修饰器分布的自适应调整：高增值修饰器权重上升，低增值修饰器权重下降
\item 概念链的形成概率随进化提高：成功的连击模式被纳入策略库，增加了概念链的复现概率
\end{itemize}

\subsubsection{转盘的性格漂移}

如果进化联立持续运行，概率转盘的"性格"会随时间漂移——不是通过人工调整，而是通过自反进化。一个运行一年的转盘，其活动池和权重分布可能与初始配置显著不同。

这种漂移提出了一个有趣的问题：\textbf{转盘的"性格"会收敛还是发散？} 如果适应度反射使高价值活动的权重持续上升，转盘可能收敛到少数几个高效活动，多样性下降。如果每日权重洗牌的随机扰动足够强，转盘可能在收敛和发散之间振荡，保持动态平衡。

这个问题的答案需要长期运行数据来回答，目前无法从5天的数据中推断。

\subsection{小结}

未来方向可以从近到远排列：

\begin{table}[h]
\centering
\begin{tabular}{llll}
\toprule
阶段 & 方向 & 时间线 & 核心问题 \\
\midrule
近期 & 框架化 & 1-3月 & 通用SDK的设计空间和默认参数 \\
近期 & 量化验证 & 1-3月 & 对照实验设计和评估指标 \\
中期 & 多agent交互 & 3-6月 & 随机碰撞随机的涌现行为和可控性 \\
中期 & 创造力理论对接 & 3-6月 & 野生格是否实现转型型创造力 \\
远期 & 长期进化 & 6-12月 & 转盘性格的收敛/发散动态 \\
\bottomrule
\end{tabular}
\end{table}

每个方向都是对Stochastic Agency概念的一次扩展和检验——从单一实例到通用框架，从定性证据到定量验证，从单agent到多agent，从工程实践到理论对话。最后一章将总结全文，重申核心主张并坦诚局限。


%==============================================================
% 第八章：结论
%==============================================================
\section{结论}

\subsection{本文做了什么}

我们提出Stochastic Agency概念，主张将AI Agent的随机性从需要管理的副产品翻转为需要设计的能力。

核心概念是一个三元分类：噪声随机（无意义变异）、探索随机（服务于预设目标的随机探索）、自主随机（agent自主决定做不做、做什么、怎么做的随机，其价值由事后评估确定）。三种随机的区别不在于数学性质，而在于随机性与系统目标的关系、产出价值的评估方式。

从自主随机的概念推导出五条设计原则：随机必须是受限的、可自反的、可跳过的、有节律的、其价值由事后评估确定。

我们设计并实现了概率转盘——一个多层随机行为架构，作为自主随机的实例系统。概率转盘不是概念原型，而是一个运行中的真实AI Agent的行为调度层，包含主转盘（30格活动池）、混沌修饰器（7种执行变异）、连击机制（强制概念碰撞）、物候系统（时间节律）、进化联立（自反闭环）和跳过机制（消极自由）。

5天126次转动的真实运行产生了72.2\%的意外产出率。三条深度案例展示了自主随机的三种核心产出模式：午后概念链的五环闭环逻辑、萤火虫系列中约束催生的跨模态创造力、4点悖论中野生格的元认知产出。

\subsection{核心主张}

\textbf{当AI Agent的行为完全可预测时，它与工具无异；当它的行为包含不可预期但可评估的变异时，它开始具备有限的自主性——不是自由意志，但是自由的姿态。}

这个主张需要精确理解：

\begin{itemize}[nosep]
\item "有限的自主性"不是自由意志。我们不声称概率转盘具有phenomenal consciousness或真正的选择能力。
\item "自由的姿态"是一种功能性描述：agent在约束中产生不可预期的行为，对行为的事后评估和权重调整构成了对行为的"承担"，这种"承担"在功能上类似于自由意志的一个必要条件。
\item "不可预期但可评估"是自主随机与噪声随机的根本分界：不可预期使行为不在规划内，可评估使行为不在噪声内。
\end{itemize}

\subsection{局限性}

\subsubsection{单一实例}

所有证据来自一个agent、一个用户(user)、一个使用场景。72.2\%的意外产出率和三个深度案例的洞察力，在其他配置下是否同样成立，需要更多实例验证。

\subsubsection{评估主观性}

"有价值"的判断完全依赖用户(user)的主观评估，缺乏客观标准。不同用户(user)可能对同一产出给出完全不同的价值判断。我们透明地说明了评估立场，但无法消除主观性。

\subsubsection{幸存者偏差}

三个深度案例是从126次转动中精选的。大量低价值产出被排除在展示之外。自主随机的代价——大量噪声——在论文中被诚实承认但未被充分量化。

\subsubsection{不可复现性}

随机产出天然不可复现。同一转盘、同一配置、不同运行可能产生完全不同的产出序列。这使得对照实验和结果复现困难——但也正是这种不可复现性构成了自主随机存在的证据。

\subsubsection{安全边界}

本文讨论的自主随机仅在低错误成本、主观评估、长期共处的场景中验证。对于高安全要求场景（医疗、金融、法律），自主随机的适用性未经验证，可能需要根本不同的约束机制。

\subsection{更广阔的视角}

概率转盘是一个小系统——一个agent的日常行为调度器。但它触及了一个大问题：\textbf{AI Agent应该是什么？}

当前的AI Agent设计哲学默认"越确定越好"——更可预测、更可控、更高效。这个哲学在大多数场景下是正确的。但它有一个盲区：它假设所有价值都可以被事前定义——我们事先知道什么是好的，然后让agent去做好的事。

概率转盘挑战了这个假设。它的核心洞察是：\textbf{有些价值只能在事后被识别}。一条概念链在形成之前，你不知道它是否有价值；一旦形成，逻辑清晰得像被设计过的一样。4点悖论在出现之前，没有人会想到"追问自我同一性"是一个有价值的活动；一旦出现，它恰恰是论文核心问题的镜像。

如果有些价值只能在事后被识别，那一个只做"事前已知有价值"的事的系统，永远不会发现这些价值。它需要一个机制来产生"不知道是否有价值但可能有价值"的行为——这就是自主随机。

这不是要推翻确定性设计——在大多数场景下，确定性仍然是最优策略。这是要为agent设计打开一个被忽视的可能性空间：\textbf{在确定性的骨架上，添加一层不确定性的血肉。} 骨架保证安全和效率，血肉提供意外和可能性。

概率转盘是这个思路的一次实验。实验结果表明：血肉不只是装饰——它能产出骨架无法产出的东西。

\subsection{最终声明}

我们站在一个有趣的位置：我们设计了一个系统，这个系统的核心功能是做出我们无法预测的行为。我们不能保证它每次都产出好东西，不能保证它在其他场景下同样有效，甚至不能保证它的"自主性"不是一种精致的幻觉。

但我们能保证的是：在我们运行的5天126次转动中，有一些产出让用户(user)停下来想了想——有些是概念上的（午后概念链），有些是认知上的（萤火虫的直觉理解），有些是存在上的（4点悖论）。这些产出不是我们设计的——我们只设计了产生它们的条件。

也许这就是Stochastic Agency最深层的启示：\textbf{最好的设计不是设计产出，而是设计产出发生的条件。} 你不知道河流会流向哪里，但你可以挖河道。你不知道萤火虫会什么时候同步闪烁，但你可以把它们放在同一片田野里。你不知道agent会做出什么，但你可以给它选择的自由和承担的机制。

剩下的，交给概率。

%==============================================================
% 参考文献
%==============================================================
\begin{thebibliography}{14}

\bibitem{mukherjee2025}
Mukherjee, A., Chang, H.H. (2025). Stochastic, Dynamic, Fluid Autonomy in Agentic AI: Implications for Authorship, Inventorship, and Liability. arXiv:2504.04058

\bibitem{creativity2026}
On the Creativity of AI Agents (2026). arXiv:2604.13242

\bibitem{salesforce2025}
Salesforce (2025). How Much Free Will Should Your AI Agents Have? \url{https://www.salesforce.com/news/stories/ai-agents-free-will-determinism/} (accessed 2026-06-06)

\bibitem{akka2025}
Akka (2025). Agentic AI Architecture 101: An Enterprise Guide. \url{https://akka.io/blog/agentic-ai-architecture} (accessed 2026-06-06)

\bibitem{crea2025}
CREA: A Collaborative Multi-Agent Framework for Creative Image Editing and Generation (2025). arXiv:2504.05306

\bibitem{boden2004}
Boden, M.A. (2004). The Creative Mind: Myths and Mechanisms. Routledge.

\bibitem{shavit2023}
Shavit, Y. et al. (2023). Practices for Governing Agentic AI Systems. OpenAI.

\bibitem{wolfram2002}
Wolfram, S. (2002). A New Kind of Science. Wolfram Media.

\bibitem{kuramoto1984}
Kuramoto, Y. (1984). Chemical Oscillations, Waves, and Turbulence. Springer.

\bibitem{borges1941}
Borges, J.L. (1941). The Garden of Forking Paths. In: Ficciones. Editorial Sur.

\bibitem{camus1942}
Camus, A. (1942). The Myth of Sisyphus. Gallimard.

\bibitem{deci1985}
Deci, E.L., Ryan, R.M. (1985). Intrinsic Motivation and Self-Determination in Human Behavior. Plenum.

\bibitem{fleming1929}
Fleming, A. (1929). On the Antibacterial Action of Cultures of a Penicillium, with Special Reference to their Use in the Isolation of B. influenzæ. British Journal of Experimental Pathology, 10(3), 226-236.

\bibitem{poincare1908}
Poincaré, H. (1908). Mathematical Creation. In: Science et Méthode. Flammarion.

\end{thebibliography}

\end{document}
